% Ad Atticum 16.2 (ad-atticum-16-02) --- English translation
% Translated by Claude (Anthropic) from the Latin (Perseus).
% Style guide: STYLE.md.

\ciceroLetterOpener{CICERO ATTICO salutem}

\ciceroSection{1} On the sixth before the Ides I received two letters, one by my own courier, the other by Brutus's. About the Buthrotian affair the report had been quite different in these parts, but with much else this too must be borne. I sent Eros back sooner than I had arranged, so that there should be someone for Hortensius and \dag for those to whom he says he had at any rate fixed the Ides\dag. Hortensius indeed is shameless. For nothing is owed to him except out of the third installment, which falls on the Kalends of Sextilis; and even of that very installment the greater part has been paid him some time in advance. But Eros will see to this on the Ides. As for Publilius, since payment has to be passed through the books, I think no delay should be made. But when you see how much of our right we have given up --- in paying down 200,000 sesterces cash out of the 400,000 still outstanding, and assigning the balance to be paid back --- you can, if it seems right to you, say to him that he ought to wait on our convenience, given that we have made so great a sacrifice of our legal right.

\ciceroSection{2} But please, my dear Atticus --- do you see how coaxing I am? --- while you are at Rome, conduct, direct, govern all my affairs in such a way that you expect nothing from me. For although what remains is well enough fitted to meet my obligations, still it often happens that those who owe do not answer to the day. If anything of the sort occurs, do not let my reputation suffer rather than anything else: by a loan if need be, by a sale even if matters force it, you shall save me.

\ciceroSection{3} Brutus was pleased with your letter. For I was with him many hours at Nesis, having received your letter a little before. He seemed to me to be delighting in the \textit{Tereus} and to feel greater gratitude to Accius than to Antony. To me, though, the gladder these things are, the more vexation and irritation I feel that the people of Rome are spending their hands not in defending the commonwealth but in applauding. Indeed, the spirits of those fellows seem to me to be set on fire to display their wickedness all the more openly. But still --- \textit{provided they suffer something, let them suffer anything at all}.

\ciceroSection{4} That my decision, as you say, is more and more praised every day --- I am not displeased to hear, and I was waiting to see whether you would write to me anything about it. For I myself was running into varied opinions. Indeed, that was why I was dragging the matter out --- so that the option might remain open as long as possible. But since we are being thrust out with a pitchfork, I am thinking of Brundisium. The legions, it seems, will be easier and surer to avoid than the pirates who are said to be in sight. Sestius was expected on the sixth before the Ides, but had not arrived, so far as I know. Cassius had arrived with his little squadron. After seeing him, I was thinking on the fifth before the Ides of moving to the Pompeianum, and then to Aeculanum. You know the rest. About Tutia, I had thought as much.

\ciceroSection{5} About Aebutius I do not believe it, nor in any case do I care more than you do. To Plancus and Oppius I have written, since you had asked --- but if it seems right to you, do not feel bound to deliver them. For when they have done everything for your sake, I am afraid they may judge my letters superfluous. To Oppius at any rate, by all means --- whom I have always known to be your warmest friend. But as you wish.

\ciceroSection{6} Since you write that you mean to winter in Epirus, you will do me a kindness if you come there before I, on your authority, am due to return to Italy. Letters to me as often as possible: on lighter matters, by anyone you can come by; but if there is anything of more weight, send it from your house. The Heracleidean treatise \textit{[Greek: H\=erakleideion]}, if we reach Brundisium safely, I shall attack. As for the \textit{De Gloria}, I have sent it to you. Guard it, then, as you usually do --- but let the excerpts \textit{[Greek: eklogai]} which Salvius is to read be marked, when, at a dinner, that is, he has secured a sympathetic audience. They please me greatly; I should rather they pleased you. Again and again, farewell.
